
%%% ВЫБОР УСТРОЙСТВА %%%

\IfStrEqCase{\devicecode}{%
{ЮТКБ.656122.611-06.02 РЭ}{%ДЗАТ
% код для этого типа
\def\fileExt{at}
\def\AbzazOne{
\item
Устройство осуществляет выдачу сигнала пуска пожаротушения автотрансформатора. Пуск пожаротушения АТ осуществляется при срабатывании дифференциальной защиты или газовой защиты АТ на отключение.
Минимальная длительность импульса выходного сигнала пуска пожаротушения определяется уставкой \textbf{«Tимп АПТ»}.
}
}%
{ЮТКБ.656122.611-11.01 РЭ}{%ШР
% код для этого типа
\def\fileExt{shr}
\def\AbzazOne{
\item
Устройство осуществляет выдачу сигнала пуска пожаротушения реактора. Пуск пожаротушения реактора осуществляется при срабатывании дифференциальной защиты или газовой защиты на отключение.
Минимальная длительность импульса выходного сигнала пуска пожаротушения определяется уставкой \textbf{«Tимп АПТ»}.
}
}%
}
[
\def\fileExt{??}
\def\AbzazOne{??}
]


\color{unimain}{\subsection{Автоматика пуска пожаротушения}}
\color{black}

\begin{enumerate}

\AbzazOne

\item
Алгоритм работы логики пуска АПТ приведен на рисунке \ref{pic:APT}. Уставки логики пуска АПТ приведены в таблице \ref{app_set:apt}.

\medskip
\begin{figure}[H]
\centering
\includegraphics[width=0.85\textwidth,height=0.85\textheight,keepaspectratio]{\fbpath/06.02 ДЗАТ/041.1141 АПТ/img/APT_\fileExt.pdf}
\captionof{figure}{Функциональная схема логики пуска АПТ}\label{pic:APT}
\end{figure}


\end{enumerate}